\documentclass[11pt]{article}

\usepackage{amsmath,amssymb,amsthm}
\usepackage{mathtools}
\usepackage[margin=1.2in]{geometry}
\usepackage{hyperref}
\usepackage{xcolor}
\usepackage{listings}
\usepackage{tikz}
\usepackage{tikz-cd}
\usepackage{mathpartir}
\usepackage{newunicodechar}
\usepackage{setspace}

\onehalfspacing

% ---- Unicode glyphs ----
\newunicodechar{→}{\ensuremath{\to}}
\newunicodechar{↦}{\ensuremath{\mapsto}}
\newunicodechar{↔}{\ensuremath{\leftrightarrow}}
\newunicodechar{≠}{\ensuremath{\neq}}
\newunicodechar{≤}{\ensuremath{\leq}}
\newunicodechar{≥}{\ensuremath{\geq}}
\newunicodechar{∀}{\ensuremath{\forall}}
\newunicodechar{∃}{\ensuremath{\exists}}
\newunicodechar{∘}{\ensuremath{\circ}}
\newunicodechar{·}{\ensuremath{\cdot}}
\newunicodechar{∈}{\ensuremath{\in}}
\newunicodechar{≃}{\ensuremath{\simeq}}
\newunicodechar{≅}{\ensuremath{\cong}}
\newunicodechar{∧}{\ensuremath{\wedge}}
\newunicodechar{∨}{\ensuremath{\vee}}
\newunicodechar{¬}{\ensuremath{\neg}}
\newunicodechar{⊢}{\ensuremath{\vdash}}
\newunicodechar{⟨}{\ensuremath{\langle}}
\newunicodechar{⟩}{\ensuremath{\rangle}}
\newunicodechar{⇑}{\ensuremath{\mathord{\uparrow}}}
\newunicodechar{∎}{\ensuremath{\blacksquare}}
\newunicodechar{ℕ}{\ensuremath{\mathbb{N}}}
\newunicodechar{ℤ}{\ensuremath{\mathbb{Z}}}
\newunicodechar{ℝ}{\ensuremath{\mathbb{R}}}
\newunicodechar{α}{\ensuremath{\alpha}}
\newunicodechar{σ}{\ensuremath{\sigma}}
\newunicodechar{φ}{\ensuremath{\varphi}}
\newunicodechar{₁}{\textsubscript{1}}
\newunicodechar{₂}{\textsubscript{2}}
\newunicodechar{≡}{\ensuremath{\equiv}}
\newunicodechar{∣}{\ensuremath{\mid}}
\newunicodechar{∑}{\ensuremath{\textstyle\sum}}
\newunicodechar{⁻}{\ensuremath{{}^{-}}}
\newunicodechar{¹}{\ensuremath{{}^{1}}}

% ---- Lean listing style ----
\definecolor{leanbg}{RGB}{247,247,250}
\definecolor{leankw}{RGB}{0,90,160}
\definecolor{leancm}{RGB}{120,120,120}
\lstdefinestyle{lean}{
  backgroundcolor=\color{leanbg},
  basicstyle=\ttfamily\small,
  keywordstyle=\color{leankw}\bfseries,
  commentstyle=\color{leancm}\itshape,
  morekeywords={theorem,lemma,def,fun,Prop,Type,variable,where,
    Field,Fintype,Finite,DecidableEq,CommRing,CommSemiring,Ring,CharP,Function,
    Injective,Surjective,Bijective,by,rfl,have,rw,exact,ext,simp,
    iterateFrobenius,RingHom,Commute,intro,convert,apply,using,
    IsAPN,Nat,Coprime,Odd,Even},
  showstringspaces=false,
  columns=fullflexible,
  extendedchars=true,
  frame=single,
  framesep=6pt,
  xleftmargin=4pt,
  aboveskip=12pt,
  belowskip=14pt,
}
\lstset{style=lean}

\newtheorem{lemma}{Lemma}

% boxed "what just happened" note
\newcommand{\note}[1]{%
  \par\smallskip
  \noindent\colorbox{leanbg}{\parbox{\dimexpr\linewidth-2\fboxsep}{\small #1}}%
  \par\smallskip}

\title{\textbf{Three Frobenius proofs, read slowly}\\[4pt]
\large Lean\,4 $\;\longleftrightarrow\;$ mathematics $\;\longleftrightarrow\;$ type theory}
\author{}
\date{}

\begin{document}
\maketitle

\begin{abstract}
\noindent
Three small proofs. Char $2$.\\[2pt]
$\bullet$\quad \texttt{frob\_bijective}: $x \mapsto x^{2^k}$ is a bijection.\\[2pt]
$\bullet$\quad \texttt{frob\_add}: $(x+y)^{2^k} = x^{2^k}+y^{2^k}$.\\[2pt]
$\bullet$\quad \texttt{sq\_add}: $(x+y)^2 = x^2+y^2$.\\[6pt]
\noindent
We read each line as: \emph{math} $\;\leftrightarrow\;$ \emph{Lean code} $\;\leftrightarrow\;$ \emph{a typing judgement} $\Gamma \vdash e : \tau$.
\end{abstract}

\bigskip

% =====================================================================
\section*{Setting}

\noindent
Fix the context $\Gamma$ once. It is the list of assumptions in scope.

\[
\Gamma \;=\;
\bigl(\, F : \mathrm{Type},\;\;
[\,\mathrm{Field}\;F\,],\;\;
[\,\mathrm{Fintype}\;F\,],\;\;
[\,\mathrm{CharP}\;F\;2\,] \,\bigr)
\]

\begin{itemize}\itemsep4pt
  \item $F$ is a \textbf{field}.
  \item $F$ is \textbf{finite} ($\mathrm{Fintype}$ $\Rightarrow$ $\mathrm{Finite}$).
  \item $F$ has \textbf{characteristic $2$}: $\;1+1 = 0$, so $\;x + x = 0$.
\end{itemize}

\note{In Lean, \texttt{[...]} are \emph{instance} arguments. They are found
automatically by type-class search. So $\Gamma$ is really
``$F$ plus three instances the elaborator carries around for free''.}

\bigskip

% =====================================================================
\section{\texttt{frob\_bijective} : the map $x \mapsto x^{2^k}$ is a bijection}

\subsection*{Statement}

\[
\forall k : \mathbb{N}.\quad
\mathrm{Bijective}\;\bigl(\, \lambda x{:}F.\; x^{2^{k}} \,\bigr)
\]

\begin{lstlisting}
lemma frob_bijective (k : ℕ) :
    Function.Bijective (fun x : F => x ^ (2 ^ k)) := by
  have : (fun x : F => x ^ (2 ^ k)) = iterateFrobenius F 2 k := by
    ext x; simp [iterateFrobenius_def]
  rw [this]
  exact Finite.injective_iff_bijective.mp (iterateFrobenius F 2 k).injective
\end{lstlisting}

As a typing goal, the proof term we must build is:
\[
\Gamma,\; k:\mathbb{N}
\;\;\vdash\;\;
?\;:\;\mathrm{Bijective}\;(\lambda x.\,x^{2^k})
\]

\bigskip
\subsection*{The plan}

The bare lambda $\lambda x.\,x^{2^k}$ is \emph{the same map} as a bundled
object Mathlib already knows: the iterated Frobenius
$\;\varphi^{k}$ where $\varphi(x)=x^{2}$.

\[
\underbrace{\lambda x.\,x^{2^{k}}}_{\text{raw function}}
\;\;=\;\;
\underbrace{\;\mathrm{iterateFrobenius}\;F\;2\;k\;}_{\text{a bundled ring hom } F \to_{+*} F}
\]

Once we know the raw map \emph{is} a ring hom on a \emph{finite} field,
two facts finish it:
\begin{enumerate}\itemsep4pt
  \item a ring hom out of a field is \textbf{injective};
  \item on a \textbf{finite} type, injective $\Rightarrow$ bijective (pigeonhole).
\end{enumerate}

\bigskip
\subsection*{Step 1 --- replace the raw map by the bundled one}

\begin{lstlisting}
have : (fun x : F => x ^ (2 ^ k)) = iterateFrobenius F 2 k := by
  ext x; simp [iterateFrobenius_def]
\end{lstlisting}

\begin{itemize}\itemsep4pt
  \item \texttt{have : ... := ...} introduces a local proof of an equality
        \emph{of functions}.
  \item \texttt{ext x} is \textbf{function extensionality}: to prove $f = g$
        it suffices to prove $\forall x,\ f\,x = g\,x$.
        \[
        \inferrule{\;\forall x.\; f\,x = g\,x\;}{\; f = g \;}\;\textsc{(funext)}
        \]
  \item \texttt{simp [iterateFrobenius\_def]} closes
        $\;x^{2^k} = (\mathrm{iterateFrobenius}\;F\;2\;k)\,x\;$
        because the rewrite rule
        \[
        \mathrm{iterateFrobenius\_def} \;:\;
        (\mathrm{iterateFrobenius}\;R\;p\;n)\,x \;=\; x^{\,p^{\,n}}
        \]
        turns the right side into $x^{2^k}$, leaving $x^{2^k}=x^{2^k}$, i.e. \texttt{rfl}.
\end{itemize}

\note{What is \texttt{iterateFrobenius F 2 k}? It is the $k$-fold composite of the
Frobenius endomorphism $\varphi:x\mapsto x^{2}$, packaged as a term of type
\texttt{F →+* F} (a \emph{ring homomorphism}). The arrow $\Uparrow$ / \texttt{⇑}
silently turns this bundle back into a plain function when needed.}

\bigskip
\subsection*{Step 2 --- rewrite the goal}

\begin{lstlisting}
rw [this]
\end{lstlisting}

\texttt{rw} substitutes left $\to$ right using the equation \texttt{this}.
The goal mutates:
\[
\mathrm{Bijective}\;(\lambda x.\,x^{2^k})
\quad\xrightarrow{\;\texttt{rw [this]}\;}\quad
\mathrm{Bijective}\;\bigl(\,\mathord{\uparrow}(\mathrm{iterateFrobenius}\;F\;2\;k)\,\bigr)
\]

\bigskip
\subsection*{Step 3 --- injective, then bijective}

\begin{lstlisting}
exact Finite.injective_iff_bijective.mp (iterateFrobenius F 2 k).injective
\end{lstlisting}

Read it right to left.

\[
\underbrace{(\mathrm{iterateFrobenius}\;F\;2\;k).\mathrm{injective}}_{:\;\mathrm{Injective}\;\varphi^{k}}
\qquad\xmapsto{\;\;\texttt{.mp}\;\;}\qquad
\underbrace{\mathrm{Bijective}\;\varphi^{k}}_{\text{goal}}
\]

\begin{itemize}\itemsep4pt
  \item \texttt{(...).injective} : a ring hom on a field is injective
        (its kernel is an ideal of a field, hence $\{0\}$).
        \[
        \Gamma \;\vdash\; \varphi^{k}.\mathrm{injective}
        \;:\; \mathrm{Injective}\;(\mathord{\uparrow}\varphi^{k})
        \]
  \item \texttt{Finite.injective\_iff\_bijective} : for a self-map of a
        \emph{finite} type,
        \[
        \mathrm{Injective}\,f \;\leftrightarrow\; \mathrm{Bijective}\,f .
        \]
        This is the pigeonhole principle:
        \[
        \inferrule
          {\;[\mathrm{Finite}\;F]\quad f : F \to F \quad \mathrm{Injective}\,f\;}
          {\;\mathrm{Surjective}\,f\;}
        \]
        and injective $\wedge$ surjective $=$ bijective.
  \item \texttt{.mp} takes the forward direction of that $\leftrightarrow$.
  \item \texttt{exact} says: \emph{this term has exactly the goal's type}. Done.
\end{itemize}

\bigskip
\subsection*{The whole derivation as one tree}

\begin{mathpar}
\inferrule*[Right=\textsc{mp}]
{
  \inferrule*[Right=\textsc{finite}]
  { [\mathrm{Finite}\;F] }
  { \mathrm{Injective}\,\varphi^k \leftrightarrow \mathrm{Bijective}\,\varphi^k }
  \\
  \inferrule*[Right=\textsc{field}]
  { \varphi^k : F \to_{+\!*} F }
  { \mathrm{Injective}\,\varphi^k }
}
{ \mathrm{Bijective}\,\varphi^k }
\end{mathpar}

\bigskip
\subsection*{Picture}

\[
\begin{tikzcd}[column sep=large, row sep=large]
F \arrow[r, "{\lambda x.\,x^{2^k}}"] \arrow[d, equal] & F \arrow[d, equal] \\
F \arrow[r, "{\mathord{\uparrow}\,\varphi^{k}}"', "{\sim}"] & F
\end{tikzcd}
\qquad
\begin{tikzcd}[row sep=tiny]
\text{inj} \arrow[r, Rightarrow, "{\text{finite}}"] & \text{surj} \\
& \text{bij}
\arrow[from=1-1, to=2-2, phantom, "{\wedge}" {yshift=2pt}]
\end{tikzcd}
\]

\medskip
\noindent
Top row $=$ raw map. Bottom row $=$ same map, bundled, and an isomorphism
($\sim$). The vertical equalities are Step\,1. $\hfill\blacksquare$

\bigskip\bigskip

% =====================================================================
\section{\texttt{frob\_add} : the Freshman's Dream, iterated}

\subsection*{Statement}

Now the context is only a \textbf{commutative ring} $R$ of char $2$
(no finiteness needed).
\[
\Gamma' = \bigl( R:\mathrm{Type},\;[\mathrm{CommRing}\;R],\;[\mathrm{CharP}\;R\;2] \bigr)
\]
\[
\forall k:\mathbb{N}.\;\forall x,y:R.\qquad
(x+y)^{2^{k}} \;=\; x^{2^{k}} + y^{2^{k}}
\]

\begin{lstlisting}
lemma frob_add (k : ℕ) (x y : R) :
    (x + y) ^ (2 ^ k) = x ^ (2 ^ k) + y ^ (2 ^ k) :=
  add_pow_char_pow_of_commute 2 k (Commute.all x y)
\end{lstlisting}

\bigskip
\subsection*{This is a term proof --- no \texttt{by}}

There is no tactic block. The body \emph{is} the proof term. We just
\textbf{apply a library lemma} to the right arguments. The skill is
matching its type to ours.

The library lemma:
\[
\mathrm{add\_pow\_char\_pow\_of\_commute} :
\quad
\Pi\,(p:\mathbb{N})\,[\mathrm{Fact}\,p.\mathrm{Prime}]\,[\mathrm{CharP}\,R\,p]
\,(n:\mathbb{N}),\;
\]
\[
\mathrm{Commute}\;x\;y \;\longrightarrow\;
(x+y)^{p^{n}} = x^{p^{n}} + y^{p^{n}}
\]

We feed it $p := 2$ and $n := k$:

\begin{center}
\begin{tabular}{r@{\;\;}c@{\;\;}l}
\textbf{slot} & & \textbf{we supply} \\[3pt]
$p$               & $\mapsto$ & \texttt{2} \\
$[\mathrm{Fact}\,p.\mathrm{Prime}]$ & $\mapsto$ & found by instance search ($2$ is prime) \\
$[\mathrm{CharP}\,R\,2]$            & $\mapsto$ & in $\Gamma'$ \\
$n$               & $\mapsto$ & \texttt{k} \\
$\mathrm{Commute}\;x\;y$ & $\mapsto$ & \texttt{Commute.all x y} \\
\end{tabular}
\end{center}

\note{\texttt{Commute.all x y} : in a \emph{commutative} ring, every pair commutes,
so $x\cdot y = y\cdot x$ is free. The general lemma only needs commutativity of
the two base elements --- here that is automatic.}

\bigskip
\subsection*{Why the theorem is true (the math under the name)}

Binomial expansion:
\[
(x+y)^{2^{k}} \;=\; \sum_{i=0}^{2^{k}} \binom{2^{k}}{i}\, x^{i} y^{2^{k}-i}.
\]
In char $2$ every interior coefficient vanishes:
\[
\binom{2^{k}}{i} \equiv 0 \pmod 2 \qquad (0 < i < 2^{k}),
\]
by Kummer / Lucas. Only the two endpoints survive:
\[
(x+y)^{2^{k}} \;=\; x^{2^{k}} + y^{2^{k}}. \tag{$\ast$}
\]
The map $\varphi:x\mapsto x^{2}$ is therefore \textbf{additive}; iterating $k$
times gives ($\ast$). Mathlib has bottled exactly this as the lemma above.

\bigskip
\subsection*{Typing judgement}

\[
\resizebox{\textwidth}{!}{$
\Gamma',\; k:\mathbb{N},\; x:R,\; y:R
\;\vdash\;
\bigl(\mathrm{add\_pow\_char\_pow\_of\_commute}\;2\;k\;(\mathrm{Commute.all}\;x\;y)\bigr)
\;:\;
(x+y)^{2^{k}} = x^{2^{k}} + y^{2^{k}}
$}
\]

The whole proof is one application node:

\begin{mathpar}
\inferrule*[Right=\textsc{app}]
{
  \mathrm{add\_pow\_char\_pow\_of\_commute}\;2\;k
  \;:\;
  \mathrm{Commute}\,x\,y \to (x+y)^{2^k}=x^{2^k}+y^{2^k}
  \\
  \mathrm{Commute.all}\;x\;y \;:\; \mathrm{Commute}\,x\,y
}
{ (x+y)^{2^k}=x^{2^k}+y^{2^k} }
\end{mathpar}

\hfill$\blacksquare$

\bigskip\bigskip

% =====================================================================
\section{\texttt{sq\_add} : one line, by specialisation}

\subsection*{Statement}

\[
\forall x,y:R.\qquad (x+y)^{2} = x^{2} + y^{2}
\]

\begin{lstlisting}
lemma sq_add (x y : R) : (x + y) ^ 2 = x ^ 2 + y ^ 2 := frob_add 1 x y
\end{lstlisting}

\bigskip
\subsection*{It is just \texttt{frob\_add} at $k = 1$}

\[
\mathrm{frob\_add}\;1\;x\;y
\;:\;
(x+y)^{2^{1}} = x^{2^{1}} + y^{2^{1}}
\]

The expected type wants the exponent $2$, the term gives $2^{1}$. These are
accepted as the same goal because they are \textbf{definitionally equal}:
\[
2^{1} \;\equiv\; 2 \cdot 2^{0} \;\equiv\; 2 \cdot 1 \;\equiv\; 2.
\]

\note{\textbf{Definitional equality} ($\equiv$) is the silent equality the kernel
checks by \emph{computation}, with no proof step. Here $2^1$ \emph{reduces} to $2$,
so Lean sees \texttt{frob\_add 1 x y} as already having the desired type. No
\texttt{rw}, no \texttt{simp} --- just reduction.}

\bigskip
\subsection*{Typing judgement}

\[
\Gamma',\; x:R,\; y:R
\;\vdash\;
\mathrm{frob\_add}\;1\;x\;y
\;:\;
\underbrace{(x+y)^{2^{1}} = x^{2^{1}} + y^{2^{1}}}_{\;\equiv\;(x+y)^2 = x^2+y^2}
\]

\begin{mathpar}
\inferrule*[Right=\textsc{conv}\;($2^1 \equiv 2$)]
{
  \mathrm{frob\_add}\;1\;x\;y \;:\; (x+y)^{2^1}=x^{2^1}+y^{2^1}
}
{ (x+y)^{2}=x^{2}+y^{2} }
\end{mathpar}

\bigskip
\subsection*{Dependency picture}

\[
\begin{tikzcd}[column sep=2.4cm, row sep=large]
\textsf{add\_pow\_char\_pow\_of\_commute}
  \arrow[r, "{p:=2,\;n:=k}"]
& \textsf{frob\_add}
  \arrow[r, "{k:=1}"]
& \textsf{sq\_add}
\end{tikzcd}
\]

\medskip
\noindent
Each arrow is ``substitute, and let the kernel reduce''. The three lemmas
form a tiny tower: a Mathlib root, one named iteration, one specialisation.
\hfill$\blacksquare$

\bigskip\bigskip

% =====================================================================
\section*{One-screen summary}

\begin{center}
\renewcommand{\arraystretch}{1.6}
\begin{tabular}{l l l}
\textbf{lemma} & \textbf{idea} & \textbf{engine} \\
\hline
\texttt{frob\_bijective} &
$x\mapsto x^{2^k}$ is $1\!\!-\!\!1$ and onto &
inj.\ (field) $+$ finite $\Rightarrow$ bij. \\
\texttt{frob\_add} &
$(x+y)^{2^k}=x^{2^k}+y^{2^k}$ &
Freshman's Dream, char $2$ \\
\texttt{sq\_add} &
$(x+y)^2=x^2+y^2$ &
\texttt{frob\_add} at $k=1$,\;\; $2^1\equiv 2$ \\
\end{tabular}
\end{center}

\bigskip
\noindent
Three flavours of proof, one per lemma:
\[
\underbrace{\texttt{frob\_bijective}}_{\text{tactic: have / rw / exact}}
\qquad
\underbrace{\texttt{frob\_add}}_{\text{term: apply a library lemma}}
\qquad
\underbrace{\texttt{sq\_add}}_{\text{term: specialise + reduce}}
\]

\bigskip\bigskip

% =====================================================================
% =====================================================================
\part*{Follow-up note: one thread, bottom to top}
\addcontentsline{toc}{part}{Follow-up note}

\bigskip
\noindent
The two char-$2$ facts above are not curiosities. They are the
\textbf{two hypotheses} of a single propagation lemma, and that lemma is
exactly the lever that lifts the Kasami APN theorem from \emph{odd $k$
only} to \emph{all $k$}.

This note follows that one thread, from the two tiny lemmas up to the
headline result. Every step is real code from
\texttt{RequestProject/Kasami/Classical/KasamiEvenK.lean}.

\bigskip
\noindent
The thread, named:
\begin{center}
\scalebox{0.74}{%
\begin{tikzcd}[ampersand replacement=\&, column sep=1.1cm, row sep=large]
\textsf{frob\_additive} \arrow[dr]
  \& \& \& \\
  \& \textsf{apn\_comp\_additive\_bijective}
      \arrow[r]
  \& \textsf{apn\_frob\_twist}
      \arrow[r]
  \& \textsf{kasami\_apn\_of\_complement}
      \arrow[d] \\
\textsf{frob\_bijective\_field} \arrow[ur]
  \& \& \& \textsf{kasami\_is\_apn\_general}
\end{tikzcd}%
}
\end{center}

\medskip
\noindent
Left: the two char-$2$ inputs. Middle: a generic preservation principle and
its Frobenius specialisation. Right: the Kasami exponent congruence feeds
that principle, and the case split finishes.

\bigskip

% ---------------------------------------------------------------------
\section{The hinge: APN is preserved by an additive bijection}

First recall what \textbf{APN} means in this project
(file \texttt{KasamiAPN.lean}):

\begin{lstlisting}
def IsAPN {F : Type*} [Field F] [CharP F 2] (f : F → F) : Prop :=
  ∀ (a : F), a ≠ 0 → ∀ (x y : F),
    f (x + a) + f x = f (y + a) + f y → y = x ∨ y = x + a
\end{lstlisting}

\note{In char $2$, $u - v = u + v$, so a ``differential collision'' is just the
equation $f(x+a)+f(x) = f(y+a)+f(y)$. \textbf{APN} $=$ ``every collision is
trivial'': $y$ must be one of the only two forced solutions $x$ or $x+a$.}

\bigskip
\noindent
The hinge lemma says: \textbf{post-composing an APN map with an additive
bijection $σ$ stays APN.}

\begin{lstlisting}
lemma apn_comp_additive_bijective {F : Type*} [Field F] [CharP F 2]
    {f : F → F} (hf : IsAPN f)
    {σ : F → F} (hσ_bij : Function.Bijective σ)
    (hσ_add : ∀ x y, σ (x + y) = σ x + σ y) :
    IsAPN (σ ∘ f) := by
  intro a ha x y hxy
  convert hf a ha x y _ using 1
  exact hσ_bij.injective ( by aesop )
\end{lstlisting}

\subsection*{Why this is the right abstraction}

Look at a collision for $σ ∘ f$ and push $σ$ through:
\[
σ(f(x{+}a)) + σ(f(x))
\;\stackrel{\text{\,$hσ\_add$\,}}{=}\;
σ\bigl(f(x{+}a)+f(x)\bigr).
\]
So a collision of $σ∘f$ is $σ$ applied to a collision of $f$:
\[
σ\bigl(f(x{+}a)+f(x)\bigr) = σ\bigl(f(y{+}a)+f(y)\bigr).
\]
Now $σ$ \textbf{injective} cancels:
\[
f(x{+}a)+f(x) = f(y{+}a)+f(y),
\]
and $f$ being APN delivers $y=x \vee y=x+a$ --- which is exactly the goal,
since $σ∘f$ has the same two forced solutions.

\note{The proof body is three lines, and each maps to one phrase above.
\texttt{intro a ha x y hxy} unpacks the collision hypothesis \texttt{hxy} for
$σ∘f$. \texttt{convert hf a ha x y \_ using 1} reduces the goal to the
remaining premise: that $f$ also collides at $x,y$. \texttt{hσ\_bij.injective
(by aesop)} discharges it --- \texttt{aesop} rewrites the $σ∘f$ collision
through \texttt{hσ\_add} into $σ(\ldots)=σ(\ldots)$, and
\texttt{.injective} strips the outer $σ$.}

\subsection*{As an inference rule}

\begin{mathpar}
\inferrule*[Right=\textsc{apn-transport}]
{
  \mathrm{IsAPN}\,f
  \\
  \mathrm{Bijective}\,σ
  \\
  \forall x\,y.\; σ(x{+}y)=σ\,x+σ\,y
}
{ \mathrm{IsAPN}\,(σ ∘ f) }
\end{mathpar}

\noindent
The two side premises are precisely the shapes of
\texttt{frob\_bijective} and \texttt{frob\_add}. That is the whole point of
this note.

\bigskip

% ---------------------------------------------------------------------
\section{Plugging in Frobenius: the twist lemma}

Now take $σ = $ Frobenius. In this file the two inputs reappear as:

\begin{lstlisting}
lemma frob_additive {F : Type*} [CommSemiring F] [CharP F 2]
    (j : ℕ) (x y : F) :
    (x + y) ^ (2 ^ j) = x ^ (2 ^ j) + y ^ (2 ^ j) :=
  add_pow_char_pow (p := 2) (n := j) x y

lemma frob_bijective_field {F : Type*} [Field F] [Fintype F] [CharP F 2]
    (j : ℕ) : Function.Bijective (fun x : F => x ^ (2 ^ j)) := by
  have : (fun x : F => x ^ (2 ^ j)) = iterateFrobenius F 2 j := by
    ext x; simp [iterateFrobenius_def]
  rw [this]
  exact Finite.injective_iff_bijective.mp (iterateFrobenius F 2 j).injective
\end{lstlisting}

\note{These are literally \texttt{frob\_add} and \texttt{frob\_bijective} from
Part\,I, re-stated in this file under the names the Kasami development uses.
\texttt{frob\_additive} is the additivity premise; \texttt{frob\_bijective\_field}
is the bijectivity premise.}

\bigskip
\noindent
Feed both into the hinge and you get: \textbf{multiplying an APN exponent by
a power of $2$ keeps it APN.}

\begin{lstlisting}
lemma apn_frob_twist {F : Type*} [Field F] [Fintype F] [CharP F 2]
    {d : ℕ} (hd : IsAPN (fun x : F => x ^ d))
    (j : ℕ) :
    IsAPN (fun x : F => x ^ (d * 2 ^ j)) := by
  have heq : (fun x : F => x ^ (d * 2 ^ j)) =
      (fun x => x ^ (2 ^ j)) ∘ (fun x => x ^ d) := by
    ext x; simp [pow_mul]
  rw [heq]
  exact apn_comp_additive_bijective hd (frob_bijective_field j) (frob_additive j)
\end{lstlisting}

\subsection*{The one line that matters}

The last line is the entire thread compressed:
\[
\underbrace{\texttt{apn\_comp\_additive\_bijective}}_{\textsc{apn-transport}}
\;
\underbrace{\texttt{hd}}_{\mathrm{IsAPN}\,(x^{d})}
\]
\[
\underbrace{\texttt{(frob\_bijective\_field j)}}_{\mathrm{Bijective}\,(x^{2^j})}
\;
\underbrace{\texttt{(frob\_additive j)}}_{σ(x+y)=σx+σy}.
\]
The rewrite \texttt{heq} uses $x^{d\cdot 2^j} = (x^{2^j})\!\circ\!(x^d)$,
i.e. \texttt{pow\_mul}, to expose the composite $σ∘f$ with
$f = x^{d}$ and $σ = x^{2^j}$.

\begin{mathpar}
\inferrule*[Right=\textsc{apn-transport}]
{
  \mathrm{IsAPN}\,(x^{d})
  \\
  \mathrm{Bijective}\,(x^{2^j})\;\textsc{(frob)}
  \\
  (x{+}y)^{2^j}=x^{2^j}{+}y^{2^j}\;\textsc{(frob)}
}
{ \mathrm{IsAPN}\,(x^{d\cdot 2^j}) }
\end{mathpar}

\bigskip

% ---------------------------------------------------------------------
\section{The arithmetic: $d_k$ \emph{is} a Frobenius twist of $d_{n-k}$}

The Kasami exponent is $d_k = 2^{2k} - 2^k + 1$:

\begin{lstlisting}
def kasamiExp (k : ℕ) : ℕ := 2 ^ (2 * k) - 2 ^ k + 1
\end{lstlisting}

\noindent
The key number fact: modulo $2^{n}-1$, the exponent $d_k$ equals
$d_{n-k}$ \emph{twisted} by $2^{2k}$.

\begin{lstlisting}
lemma kasami_exp_congr_mod {k n : ℕ} (hk : 0 < k) (hkn : k < n) :
    kasamiExp k % (2 ^ n - 1) =
    (kasamiExp (n - k) * 2 ^ (2 * k)) % (2 ^ n - 1) := by ...
\end{lstlisting}

\[
d_k \;\equiv\; d_{n-k}\cdot 2^{2k} \pmod{2^{n}-1}.
\]

\note{Why mod $2^{n}-1$? On $\mathrm{GF}(2^{n})^{\times}$ every element has order
dividing $2^{n}-1$, so $x^{e}$ depends only on $e \bmod (2^{n}-1)$.
A congruence of exponents becomes an \emph{equality of power maps}.}

\bigskip
\noindent
Transported to the field, the congruence becomes a pointwise identity
(the $x=0$ case is separate, the $x\ne 0$ case uses
$x^{2^{n}}=x$ and the congruence):

\begin{lstlisting}
lemma kasami_pow_frob_identity {F : Type*} [Field F] [Fintype F] [CharP F 2]
    {n : ℕ} (hn : Fintype.card F = 2 ^ n)
    (k : ℕ) (hk : 0 < k) (hkn : k < n) (x : F) :
    x ^ kasamiExp k = (x ^ kasamiExp (n - k)) ^ (2 ^ (2 * k)) := by ...
\end{lstlisting}

\[
x^{d_k} \;=\; \bigl(x^{\,d_{n-k}}\bigr)^{2^{2k}}
\;=\;
σ\bigl(x^{\,d_{n-k}}\bigr),\qquad σ = \text{Frob}_{2k}.
\]

\bigskip

% ---------------------------------------------------------------------
\section{Closing the loop: APN of the complement}

The two halves --- ``$d_k$ is $d_{n-k}$ twisted'' and ``twisting preserves
APN'' --- snap together:

\begin{lstlisting}
lemma kasami_apn_of_complement {F : Type*} [Field F] [Fintype F] [CharP F 2]
    {n : ℕ} (hn : Fintype.card F = 2 ^ n)
    (k : ℕ) (hk : 0 < k) (hkn : k < n)
    (hapn : IsAPN (fun x : F => x ^ kasamiExp (n - k))) :
    IsAPN (fun x : F => x ^ kasamiExp k) := by
  have heq : (fun x : F => x ^ kasamiExp k) =
      (fun x : F => x ^ (kasamiExp (n - k) * 2 ^ (2 * k))) := by
    ext x; rw [pow_mul]; exact kasami_pow_frob_identity hn k hk hkn x
  rw [heq]
  exact apn_frob_twist hapn (2 * k)
\end{lstlisting}

\note{Read the last two lines as the thread closing: \texttt{heq} rewrites
$x^{d_k}$ into $x^{d_{n-k}\cdot 2^{2k}}$ (the arithmetic of \S3), then
\texttt{apn\_frob\_twist hapn (2*k)} applies the twist lemma of \S2 ---
which is the hinge of \S1 fed with the two char-$2$ lemmas.}

\begin{mathpar}
\inferrule*[Right=\textsc{complement}]
{
  \mathrm{IsAPN}\,(x^{d_{n-k}})
  \\
  x^{d_k} = \bigl(x^{d_{n-k}}\bigr)^{2^{2k}}\;\textsc{(\S3)}
}
{ \mathrm{IsAPN}\,(x^{d_k}) }
\end{mathpar}

\bigskip

% ---------------------------------------------------------------------
\section{The headline: the parity restriction disappears}

The \emph{original} Kasami theorem (\texttt{KasamiAPN.lean}) carries
\texttt{(hk\_odd : Odd k)}:

\begin{lstlisting}
theorem kasami_is_apn {F : Type*} [Field F] [Fintype F] [CharP F 2]
    {n : ℕ} (hn : Fintype.card F = 2 ^ n) (k : ℕ)
    (hk : 1 < k) (hk_odd : Odd k) (hkn : k < n)
    (hn_odd : Odd n) (hcop : Nat.Coprime k n) :
    IsAPN (fun (x : F) => x ^ (kasamiExp k)) := ...
\end{lstlisting}

\noindent
When $k$ is \textbf{even} and $n$ is \textbf{odd}, the complement
$n-k$ is \textbf{odd} with the \emph{same} coprimality --- so
\texttt{kasami\_is\_apn} applies to $n-k$, and
\texttt{kasami\_apn\_of\_complement} carries it back to $k$:

\begin{lstlisting}
theorem kasami_is_apn_even_k {F : Type*} [Field F] [Fintype F] [CharP F 2]
    {n : ℕ} (hn : Fintype.card F = 2 ^ n) (k : ℕ)
    (hk : 1 < k) (hk_even : Even k) (hkn : k < n)
    (hn_odd : Odd n) (hcop : Nat.Coprime k n)
    (hnk : n - k ≥ 2) :
    IsAPN (fun (x : F) => x ^ kasamiExp k) := by
  apply kasami_apn_of_complement hn k (by omega) hkn
  have hcop' : Nat.Coprime (n - k) n := coprime_sub_self (le_of_lt hkn) hcop
  have hodd' : Odd (n - k) := odd_sub_even hn_odd hk_even (le_of_lt hkn)
  exact kasami_is_apn hn (n - k) hnk hodd' (by omega) hn_odd hcop'
\end{lstlisting}

\note{The edge case $n-k=1$ is handled separately
(\texttt{kasami\_is\_apn\_even\_k\_edge}): there the complement exponent is the
Gold exponent $d_1 = 3 = 2^1+1$, and Gold APN is proved directly
(\texttt{gold\_is\_apn}). Same complement trick, different base case.}

\bigskip
\noindent
The unified statement --- \textbf{no parity hypothesis on $k$}:

\begin{lstlisting}
theorem kasami_is_apn_general {F : Type*} [Field F] [Fintype F] [CharP F 2]
    {n : ℕ} (hn : Fintype.card F = 2 ^ n) (k : ℕ)
    (hk : 1 < k) (hkn : k < n)
    (hn_odd : Odd n) (hcop : Nat.Coprime k n) :
    IsAPN (fun (x : F) => x ^ kasamiExp k) := by
  by_cases hk_odd : Odd k
  · exact kasami_is_apn hn k hk hk_odd hkn hn_odd hcop
  · have hk_even : Even k := Nat.not_odd_iff_even.mp hk_odd
    by_cases hnk : n - k ≥ 2
    · exact kasami_is_apn_even_k hn k hk hk_even hkn hn_odd hcop hnk
    · have hnk1 : n - k = 1 := by omega
      exact kasami_is_apn_even_k_edge hn k hk hk_even hkn hn_odd hcop hnk1
\end{lstlisting}

\note{Compare the two signatures: \texttt{kasami\_is\_apn} has
\texttt{(hk\_odd : Odd k)}; \texttt{kasami\_is\_apn\_general} does \emph{not}.
That single deleted hypothesis is what the whole thread bought.}

\bigskip

% ---------------------------------------------------------------------
\section*{The thread, on one screen}

\begin{center}
\renewcommand{\arraystretch}{1.55}
\begin{tabular}{l l}
\textbf{step} & \textbf{what it adds} \\
\hline
\texttt{frob\_additive} & $σ$ is additive \;(char $2$) \\
\texttt{frob\_bijective\_field} & $σ$ is a bijection \;(field $+$ finite) \\
\texttt{apn\_comp\_additive\_bijective} & inj.\ $+$ additive $\Rightarrow$ APN survives $σ∘(-)$ \\
\texttt{apn\_frob\_twist} & specialise $σ = x^{2^j}$: $x^{d}$ APN $\Rightarrow$ $x^{d\cdot 2^j}$ APN \\
\texttt{kasami\_exp\_congr\_mod} & $d_k \equiv d_{n-k}\cdot 2^{2k} \pmod{2^n-1}$ \\
\texttt{kasami\_pow\_frob\_identity} & $x^{d_k} = (x^{d_{n-k}})^{2^{2k}}$ on the field \\
\texttt{kasami\_apn\_of\_complement} & APN at $n-k$ $\Rightarrow$ APN at $k$ \\
\texttt{kasami\_is\_apn\_general} & drop \texttt{Odd k} \\
\end{tabular}
\end{center}

\bigskip
\noindent
One sentence: \;\emph{Frobenius is an additive bijection
(\texttt{frob\_add} $+$ \texttt{frob\_bijective}), an additive bijection
preserves APN, the Kasami exponent $d_k$ is a Frobenius twist of its
complement $d_{n-k}$, and so the odd-$k$ theorem covers the even-$k$ case by
swapping $k \leftrightarrow n-k$ --- erasing the parity restriction.}

\[
\boxed{\;
\begin{aligned}
&\texttt{frob\_add} + \texttt{frob\_bijective}
\;\Longrightarrow\; \text{APN preserved by twist}\\
&\hphantom{\texttt{frob\_add} + \texttt{frob\_bijective}\;}
\Longrightarrow\; \texttt{kasami\_is\_apn\_general}
\end{aligned}
\;}
\]

\end{document}
